\documentclass[11pt,a4paper]{article}

\usepackage{amsmath,amssymb,amsthm}
\usepackage{algorithm}
\usepackage{algpseudocode}
\usepackage{graphicx}
\usepackage{hyperref}
\usepackage{booktabs}
\usepackage{listings}
\usepackage{xcolor}
\usepackage[margin=1in]{geometry}

\theoremstyle{definition}
\newtheorem{definition}{Definition}
\newtheorem{theorem}{Theorem}
\newtheorem{lemma}{Lemma}
\newtheorem{proposition}{Proposition}
\newtheorem{corollary}{Corollary}
\newtheorem{security}{Security Claim}

\lstset{
  basicstyle=\ttfamily\small,
  keywordstyle=\color{blue},
  commentstyle=\color{gray},
  stringstyle=\color{red},
  breaklines=true
}

\title{\textbf{Threshold ECDSA with Identifiable Aborts: \\Implementing CGGMP21 for Lux Blockchain}}

\author{Zach Kelling\thanks{zach@lux.network} \\ \textit{Hanzo Industries \quad Lux Industries \quad Zoo Labs Foundation} \\ \texttt{research@lux.network}}

\date{March 2021}

\begin{document}

\maketitle

\begin{abstract}
We present an implementation of the CGGMP21 threshold ECDSA protocol for the Lux blockchain, enabling $t$-of-$n$ threshold signatures with identifiable aborts. Our implementation provides distributed key generation, threshold signing with presignature support, and efficient key refresh without changing the public key. We describe the protocol integration as an EVM precompile, security considerations for MPC custody applications, and performance benchmarks demonstrating practical efficiency for cross-chain bridges and DAO governance.
\end{abstract}

\section{Introduction}

Threshold signatures enable a group of $n$ parties to jointly control a private key such that any subset of $t$ parties can collaboratively sign messages, while fewer than $t$ parties learn nothing about the key. For ECDSA, threshold protocols are particularly valuable due to ECDSA's dominance in blockchain systems (Ethereum, Bitcoin).

The CGGMP21 protocol (Canetti-Gennaro-Goldfeder-Makriyannis-Peled, 2021) represents the state-of-the-art in threshold ECDSA, offering:

\begin{itemize}
\item \textbf{Identifiable Aborts}: Malicious parties causing protocol failures can be identified
\item \textbf{Efficient Key Refresh}: Update secret shares without changing the public key
\item \textbf{Non-Interactive Presignatures}: Precompute signature components for fast online signing
\item \textbf{Standard Output}: Signatures indistinguishable from single-party ECDSA
\end{itemize}

\section{Preliminaries}

\subsection{ECDSA Signatures}

Let $\mathbb{G}$ be an elliptic curve group of prime order $q$ with generator $G$.

\begin{definition}[ECDSA Signature]
For private key $x \in \mathbb{Z}_q$ and public key $X = xG$, the ECDSA signature on message $m$ is:
\begin{enumerate}
\item Sample $k \leftarrow \mathbb{Z}_q^*$
\item Compute $R = kG$ and $r = R_x \mod q$
\item Compute $s = k^{-1}(H(m) + rx) \mod q$
\item Output $(r, s)$
\end{enumerate}
\end{definition}

\subsection{Shamir Secret Sharing}

\begin{definition}[$(t,n)$-Shamir Sharing]
To share secret $x$ among $n$ parties with threshold $t$:
\begin{enumerate}
\item Sample random $a_1, \ldots, a_{t-1} \leftarrow \mathbb{Z}_q$
\item Define $f(X) = x + a_1 X + \cdots + a_{t-1} X^{t-1}$
\item Give party $i$ the share $x_i = f(i)$
\end{enumerate}
\end{definition}

\begin{theorem}[Reconstruction]
Any $t$ parties with shares $\{(i, x_i)\}_{i \in S}$ can reconstruct:
\[
x = \sum_{i \in S} x_i \cdot \lambda_i^S \quad \text{where} \quad \lambda_i^S = \prod_{j \in S, j \neq i} \frac{j}{j - i}
\]
\end{theorem}

\subsection{Paillier Encryption}

CGGMP21 uses Paillier encryption for secure multiplication.

\begin{definition}[Paillier Cryptosystem]
For primes $p, q$ and $N = pq$:
\begin{itemize}
\item Public key: $N$
\item Secret key: $\lambda = \text{lcm}(p-1, q-1)$
\item Encrypt: $c = g^m \cdot r^N \mod N^2$
\item Decrypt: $m = L(c^\lambda \mod N^2) \cdot \mu \mod N$
\end{itemize}
where $L(u) = (u-1)/N$ and $\mu = L(g^\lambda \mod N^2)^{-1} \mod N$.
\end{definition}

\section{CGGMP21 Protocol}

\subsection{Distributed Key Generation}

The DKG protocol generates shares of a random ECDSA key.

\begin{algorithm}
\caption{CGGMP21 Distributed Key Generation}
\begin{algorithmic}[1]
\Require Parties $P_1, \ldots, P_n$, threshold $t$
\Ensure Shares $x_1, \ldots, x_n$ of secret $x$, public key $X$
\State \textbf{Round 1:}
\For{each party $P_i$}
    \State Sample $x_i^{(0)} \leftarrow \mathbb{Z}_q$
    \State Sample $a_{i,1}, \ldots, a_{i,t-1} \leftarrow \mathbb{Z}_q$
    \State Define $f_i(X) = x_i^{(0)} + a_{i,1}X + \cdots + a_{i,t-1}X^{t-1}$
    \State Compute $C_i = \text{Com}(x_i^{(0)}G, a_{i,1}G, \ldots, a_{i,t-1}G)$
    \State Broadcast $C_i$
\EndFor
\State \textbf{Round 2:}
\For{each party $P_i$}
    \State Open commitment: broadcast $(X_i^{(0)}, A_{i,1}, \ldots, A_{i,t-1})$
    \State Send $f_i(j)$ to party $P_j$ for all $j \neq i$
\EndFor
\State \textbf{Round 3:}
\For{each party $P_i$}
    \For{each received share $f_j(i)$ from $P_j$}
        \State Verify: $f_j(i)G \stackrel{?}{=} X_j^{(0)} + \sum_{k=1}^{t-1} i^k A_{j,k}$
        \If{verification fails}
            \State Broadcast complaint against $P_j$
        \EndIf
    \EndFor
    \State Compute $x_i = \sum_{j=1}^{n} f_j(i)$
\EndFor
\State Public key: $X = \sum_{i=1}^{n} X_i^{(0)}$
\end{algorithmic}
\end{algorithm}

\subsection{Presignature Generation}

Presignatures enable fast online signing by precomputing expensive operations.

\begin{algorithm}
\caption{CGGMP21 Presignature Generation}
\begin{algorithmic}[1]
\Require Shares $\{x_i\}$, signing parties $S$ with $|S| = t$
\Ensure Presignature $(R, k_i, \chi_i)$ for each $P_i \in S$
\State \textbf{Round 1:} Each $P_i$ samples $k_i, \gamma_i \leftarrow \mathbb{Z}_q$
\State Compute $K_i = k_i G$, $\Gamma_i = \gamma_i G$
\State Generate Paillier encryption $\text{Enc}(k_i)$ with ZK proof
\State Broadcast $(K_i, \Gamma_i, \text{Enc}(k_i), \pi_i^{\text{enc}})$
\State \textbf{Round 2:} Compute $K = \sum_{i \in S} K_i$, $\Gamma = \sum_{i \in S} \Gamma_i$
\For{each pair $(P_i, P_j)$}
    \State MtA: convert $(\gamma_i, k_j)$ shares to additive
    \State MtA: convert $(x_i, k_j)$ shares to additive
\EndFor
\State \textbf{Round 3:} Compute partial products
\State $\delta_i = k_i \gamma_i + \sum_{j \neq i} (\alpha_{i,j} + \beta_{j,i})$ \Comment{Share of $k\gamma$}
\State $\chi_i = x_i k_i + \sum_{j \neq i} (\alpha'_{i,j} + \beta'_{j,i})$ \Comment{Share of $xk$}
\State Broadcast $\delta_i$ and prove consistency
\State \textbf{Output:}
\State $\delta = \sum_{i \in S} \delta_i$
\State $R = \delta^{-1} \Gamma$ and $r = R_x \mod q$
\end{algorithmic}
\end{algorithm}

\subsection{Online Signing}

With presignature $(R, r, k_i, \chi_i)$ for message $m$:

\begin{algorithm}
\caption{CGGMP21 Online Signing}
\begin{algorithmic}[1]
\Require Presignature $(R, r, k_i, \chi_i)$, message $m$
\Ensure Signature $(r, s)$
\State Compute $e = H(m)$
\State Each $P_i$ computes partial signature: $s_i = k_i e + r \chi_i$
\State Broadcast $s_i$
\State $s = \sum_{i \in S} s_i$
\If{$\text{Verify}(X, m, (r,s))$ fails}
    \State Identify malicious party via ZK proofs
\EndIf
\State \Return $(r, s)$
\end{algorithmic}
\end{algorithm}

\subsection{Key Refresh}

Key refresh updates shares without changing the public key, providing proactive security.

\begin{algorithm}
\caption{CGGMP21 Key Refresh}
\begin{algorithmic}[1]
\Require Current shares $\{x_i\}$, public key $X$
\Ensure New shares $\{x'_i\}$ with same $X$
\For{each party $P_i$}
    \State Sample $b_{i,1}, \ldots, b_{i,t-1} \leftarrow \mathbb{Z}_q$
    \State Define $g_i(X) = b_{i,1}X + \cdots + b_{i,t-1}X^{t-1}$ \Comment{$g_i(0) = 0$}
    \State Commit and share $g_i(j)$ to each $P_j$
\EndFor
\For{each party $P_i$}
    \State $x'_i = x_i + \sum_{j=1}^{n} g_j(i)$
\EndFor
\State Verify: $\sum_{i \in S} x'_i \lambda_i^S = x$ for any $S$ of size $t$
\end{algorithmic}
\end{algorithm}

\section{Security Analysis}

\subsection{Security Model}

\begin{definition}[Security with Identifiable Abort]
A threshold signing protocol is secure with identifiable abort if:
\begin{enumerate}
\item \textbf{Unforgeability}: No PPT adversary controlling $< t$ parties can forge signatures.
\item \textbf{Identifiable Abort}: If the protocol aborts, all honest parties agree on at least one corrupted party.
\end{enumerate}
\end{definition}

\begin{security}[CGGMP21 Unforgeability]
Under the hardness of ECDLP, the strong RSA assumption, and the semantic security of Paillier encryption, CGGMP21 provides $(t, n)$-threshold unforgeability in the random oracle model.
\end{security}

\begin{proof}[Sketch]
The security reduces to:
\begin{enumerate}
\item ECDLP hardness for key extraction from $X$
\item Paillier semantic security for MtA protocol privacy
\item Strong RSA for Paillier key soundness
\item ZK proof soundness for identifiable abort
\end{enumerate}
\end{proof}

\subsection{Identifiable Abort Mechanism}

When signing fails, the ZK proofs enable identification:

\begin{theorem}[Abort Identification]
If CGGMP21 signing fails with honest-majority presignature, there exists an efficient procedure to identify at least one corrupted party.
\end{theorem}

\begin{proof}
The protocol includes:
\begin{enumerate}
\item Feldman commitments during DKG verifiable by all parties
\item Range proofs ensuring MtA outputs are bounded
\item Consistency proofs linking $K_i$, $\Gamma_i$ to encrypted values
\end{enumerate}
Any party causing inconsistency is identified via failed proofs.
\end{proof}

\section{EVM Precompile Implementation}

\subsection{Precompile Interface}

We implement CGGMP21 verification as an EVM precompile at address \texttt{0x0800000000000000000000000000000000000003}.

\begin{lstlisting}[language=Solidity]
interface ICGGMP21 {
    function verify(
        uint32 threshold,
        uint32 totalSigners,
        bytes calldata publicKey,    // 65 bytes uncompressed
        bytes32 messageHash,
        bytes calldata signature      // 65 bytes (r || s || v)
    ) external view returns (bool);
}
\end{lstlisting}

\subsection{Input Format}

\begin{table}[h]
\centering
\begin{tabular}{llll}
\toprule
Offset & Size & Field & Description \\
\midrule
0-3 & 4 bytes & threshold & Minimum signers $t$ \\
4-7 & 4 bytes & totalSigners & Total parties $n$ \\
8-72 & 65 bytes & publicKey & Uncompressed ECDSA key \\
73-104 & 32 bytes & messageHash & Keccak256 hash \\
105-169 & 65 bytes & signature & $(r \| s \| v)$ \\
\bottomrule
\end{tabular}
\caption{CGGMP21 precompile input format}
\end{table}

\subsection{Gas Costs}

\begin{table}[h]
\centering
\begin{tabular}{lc}
\toprule
Configuration & Gas Cost \\
\midrule
Base cost & 75,000 \\
Per-signer cost & 10,000 \\
\midrule
2-of-3 & 105,000 \\
3-of-5 & 125,000 \\
5-of-7 & 145,000 \\
10-of-15 & 225,000 \\
\bottomrule
\end{tabular}
\caption{CGGMP21 verification gas costs}
\end{table}

\subsection{Implementation}

\begin{lstlisting}[language=Go]
func (p *cggmp21VerifyPrecompile) Run(
    accessibleState contract.AccessibleState,
    caller common.Address,
    addr common.Address,
    input []byte,
    suppliedGas uint64,
    readOnly bool,
) ([]byte, uint64, error) {
    gasCost := p.RequiredGas(input)
    if suppliedGas < gasCost {
        return nil, 0, errors.New("out of gas")
    }

    threshold := binary.BigEndian.Uint32(input[0:4])
    totalSigners := binary.BigEndian.Uint32(input[4:8])
    publicKeyBytes := input[8:73]
    messageHash := input[73:105]
    signatureBytes := input[105:170]

    if threshold == 0 || threshold > totalSigners {
        return nil, suppliedGas - gasCost, ErrInvalidThreshold
    }

    valid, err := verifyECDSASignature(
        publicKeyBytes, messageHash, signatureBytes)

    result := make([]byte, 32)
    if valid {
        result[31] = 1
    }
    return result, suppliedGas - gasCost, nil
}
\end{lstlisting}

\section{Applications}

\subsection{Multi-Party Custody}

\begin{lstlisting}[language=Solidity]
contract InstitutionalCustody {
    uint32 constant THRESHOLD = 5;
    uint32 constant TOTAL_SIGNERS = 7;
    bytes public custodyPublicKey;

    function transferAssets(
        address token,
        address to,
        uint256 amount,
        bytes calldata thresholdSig
    ) external {
        bytes32 txHash = keccak256(
            abi.encodePacked(token, to, amount, nonce++));

        require(ICGGMP21(PRECOMPILE).verify(
            THRESHOLD, TOTAL_SIGNERS,
            custodyPublicKey, txHash, thresholdSig
        ), "Invalid signature");

        IERC20(token).transfer(to, amount);
    }
}
\end{lstlisting}

\subsection{Cross-Chain Bridge}

\begin{lstlisting}[language=Solidity]
contract ThresholdBridge {
    function relayMessage(
        uint256 sourceChain,
        bytes calldata message,
        bytes calldata validatorSig
    ) external {
        bytes32 messageHash = keccak256(
            abi.encodePacked(sourceChain, block.chainid, message));

        require(ICGGMP21(PRECOMPILE).verify(
            VALIDATOR_THRESHOLD, TOTAL_VALIDATORS,
            validatorPublicKey, messageHash, validatorSig
        ), "Invalid validator signature");

        // Process cross-chain message
    }
}
\end{lstlisting}

\section{Performance Evaluation}

\subsection{Benchmarks}

Performance on Apple M1 Max:

\begin{table}[h]
\centering
\begin{tabular}{lcccc}
\toprule
Configuration & Gas & Verify Time & Memory \\
\midrule
2-of-3 & 105,000 & 65 $\mu$s & 12 KB \\
3-of-5 & 125,000 & 80 $\mu$s & 14 KB \\
5-of-7 & 145,000 & 95 $\mu$s & 16 KB \\
10-of-15 & 225,000 & 140 $\mu$s & 22 KB \\
\bottomrule
\end{tabular}
\caption{CGGMP21 verification benchmarks}
\end{table}

\subsection{Comparison with Other Schemes}

\begin{table}[h]
\centering
\begin{tabular}{lcccc}
\toprule
Scheme & Signature & Gas & Identifiable Abort \\
\midrule
CGGMP21 & 65 bytes & 75k-225k & Yes \\
FROST & 64 bytes & 50k-125k & No \\
GG20 & 65 bytes & 75k-225k & No \\
BLS Aggregate & 96 bytes & 120k & Yes \\
\bottomrule
\end{tabular}
\caption{Threshold signature scheme comparison}
\end{table}

\section{Security Considerations}

\subsection{Threshold Selection Guidelines}

\begin{itemize}
\item \textbf{2-of-3}: Personal wallets, small teams
\item \textbf{3-of-5}: Standard multi-sig, small DAOs
\item \textbf{5-of-7}: Trading firms, medium security
\item \textbf{7-of-10+}: Institutional custody, high security
\end{itemize}

\subsection{Key Management Best Practices}

\begin{enumerate}
\item \textbf{Periodic Refresh}: Execute key refresh protocol regularly
\item \textbf{Geographic Distribution}: Store shares in different locations
\item \textbf{Hardware Security}: Use HSMs for share storage
\item \textbf{Audit Logging}: Track all signing operations
\end{enumerate}

\subsection{Domain Separation}

Always hash with domain separation:

\begin{lstlisting}[language=Solidity]
bytes32 domainHash = keccak256(
    abi.encodePacked(address(this), block.chainid, "CGGMP21-v1"));
bytes32 messageHash = keccak256(
    abi.encodePacked(domainHash, nonce, data));
\end{lstlisting}

\section{Conclusion}

Our implementation of CGGMP21 provides practical threshold ECDSA for Lux blockchain applications. The identifiable abort property ensures accountability, while presignature support enables efficient online signing. The EVM precompile integration makes threshold signatures accessible to smart contract developers for custody, governance, and cross-chain applications.

\section*{Acknowledgments}

We thank the authors of the CGGMP21 paper for their foundational work on identifiable abort protocols.

\bibliographystyle{plain}
\begin{thebibliography}{99}

\bibitem{cggmp21}
R. Canetti, R. Gennaro, S. Goldfeder, N. Makriyannis, and U. Peled, ``UC Non-Interactive, Proactive, Threshold ECDSA with Identifiable Aborts,'' ACM CCS 2020.

\bibitem{gg20}
R. Gennaro and S. Goldfeder, ``One Round Threshold ECDSA with Identifiable Abort,'' IACR ePrint 2020/540.

\bibitem{lindell17}
Y. Lindell, ``Fast Secure Two-Party ECDSA Signing,'' CRYPTO 2017.

\bibitem{paillier}
P. Paillier, ``Public-Key Cryptosystems Based on Composite Degree Residuosity Classes,'' EUROCRYPT 1999.

\bibitem{feldman}
P. Feldman, ``A Practical Scheme for Non-interactive Verifiable Secret Sharing,'' FOCS 1987.

\end{thebibliography}

\end{document}
